\RequirePackage[2020-02-02]{latexrelease}  
\documentclass[12pt]{zHenriquesLab-StyleBioRxiv}
\usepackage{relsize, multirow, xcolor,xargs, placeins}
\usepackage{amsfonts, amsmath, amsthm, amssymb, siunitx} % For math fonts, symbols and environments
\usepackage{environ, etoolbox}
\usepackage{lineno} 
\definecolor{sunyGray}{RGB}{131,134,135} 
\definecolor{sunyBlue}{RGB}{0,76,147}     
\newcommandx{\cyansub}[1]{\large\textcolor{sunyBlue} \noindent \textcolor{sunyBlue} {\noindent\textbf{#1}}}                             
\newcommand{\myemph}[1]{\textbf{\textit{#1}}}
\newcommand{\Ga}{\mbox{Ga}} 
\newcommand{\N}{\mathcal{N}} 
\newcommand{\E}{\mbox{E}}                                                                                                                                                           
\newcommand\fref[1]{\normalsize{\textbf{\textit{Fig.\,\ref{#1}}}}\normalsize{}}                                                                                                    
\newcommand\dref[2]{\normalsize{\textbf{\textit{Fig.\,\ref{#1}#2}}}\normalsize{}}                                                                                                  
\newcommand\sref[1]{\normalsize{\textbf{\textit{Supplementary Fig.\,\ref{#1}}}}\normalsize{}}                                                                                       
\newcommand\stref[1]{\normalsize{\textbf{\textit{Supplementary Table\,\ref{#1}}}}\normalsize{}}                                                                                       

\usepackage{float} 
\newfloat{schematic}{tbph}{loc}
\floatname{schematic}{Schematic}
\captionsetup[schematic]{position=bottom}
\crefname{schematic}{\textit{schematic}}{charts}
\Crefname{schematic}{\textit{Schematic}}{\textit{Charts}}


\newfloat{pubdata}{tbph}{loc}
\floatname{pubdata}{\scriptsize Fig}
\crefname{pubdata}{figure}{charts}
\Crefname{pubdata}{\textit{Fig}}{\textit{Charts}}



\providecommand{\figdir}{.}
\graphicspath{{\figdir/}}


%\newfloat{mydata}{tbph}{loc}
%\floatname{mydata}{\scriptsize Fig}
%\crefname{mydata}{figure}{charts}
%\Crefname{mydata}{\textit{Fig}}{\textit{Charts}}


%\newfloat{graph2}{tbph}{lom}[chapter]
%\restylefloat*{graph2}
%\floatstyle{plain}
%\floatname{grap2}{Graph2}
%\captionsetup[graph2]{position=top}
%\newcommand{\listofGraphs}{\listof{Graph2}{List of Graphs}}
%\newsubfloat[position=bottom,listofformat=subsimple]{graph2}




\setlength\parindent{0.50cm}
\setlength{\parindent}{0cm}
\setlength{\parskip}{0.4\baselineskip} 

\leadauthor{Souaiaia} 
\begin{document}
\title{Striking Departures from Polygenic Architecture in the Tails of Complex Traits}
\shorttitle{Tails}
\onecolumn 
%\maketitle
\vspace*{5mm}
%\begin{corrauthor} \texttt{\small tade.souaiaia\at downstate.edu, clive.hoggart\at mssm.edu, paul.oreilly\at mssm.edu} \end{corrauthor}
%\begin{corrauthor} \texttt{\small tade.souaiaia\at downstate.edu, paul.oreilly\at mssm.edu} \end{corrauthor}
%\vspace*{-2mm}
%\linenumbers

\section*{Supplemental Data} 
\noindent\textcolor{sunyGray}{\rule{17.25cm}{1mm}}

See the Excel/CSV Data Tables: 
\begin{itemize} 
    \item table-commonTails.csv  
    \item table-rareSnps.csv 
    \item table-traitSummaries.csv
    \item table-glossary.csv    
    \item table-rareTails.csv
\end{itemize} 



\clearpage \newpage


\renewcommand{\figurename}{Supplemental Figure}
\setcounter{figure}{0}   

\section*{Supplemental Figures} \label{sec:supplement}
\noindent\textcolor{sunyGray}{\rule{17.25cm}{1mm}}


\FloatBarrier \begin{figure}[!h] 
\includegraphics[trim=0cm 0cm 0cm 0cm, clip, width=\linewidth]{sup1} 
\caption{{{\bf Replication of Sanjak et al. lifetime reproductive success results.}
Comparison of results among overlapping traits (without sex differences) shows replication (i.e. statistically 
significant parameters included in our model) for 19 of 20 parameters. The only exception is 
\textit{Body Fat}, in which Sanjak et al reported a positive stabilising result in males and 
females, while we selected a neutral stabilising model in the pooled sample. 
}} \label{sup1} \end{figure}
\FloatBarrier
\clearpage \newpage

\FloatBarrier \begin{figure}[!h] 
\includegraphics[trim=0cm 0cm 0cm 0cm, clip, width=\linewidth]{sup2} 
\caption{{\bf BIC model selection sensitivity analysis.}
\textbf{a,} For 17 of 74 traits, the top model selected using BIC was not significantly better fitting than 
the second best model. 
\textbf{b,} Association between BIC selected models of selection and POPout effects (top models). 
\textbf{c,} Association between BIC selected models of selection and POPout effects after instead 
using the second best model for the 17 traits in which the top model was not significantly better fitting 
than the second best model. 
} \label{sup2} \end{figure}
\FloatBarrier









\clearpage \newpage
\end{document}







